\documentclass[a4paper,12pt]{article}
\usepackage[T2A]{fontenc}
\usepackage[utf8]{inputenc}
\usepackage[russian]{babel}
\usepackage[margin=25mm]{geometry}
\usepackage{xcolor,tikz,booktabs}
\definecolor{accent}{HTML}{008080}
\definecolor{darkaccent}{HTML}{004D4D}
\definecolor{textgray}{HTML}{333333}
\pagestyle{empty}
\begin{document}
\begin{center}
\vspace*{23mm}
{\Large\color{darkaccent} МАТЕМАТИКА\par}
\vspace{15mm}
{\huge Ревью домашней работы\par}
\vspace{11mm}
{\LARGE\color{accent} Цилиндр: площадь и объём\par}
\vspace{19mm}
{\large Николь Саркисьянц\par}
\vspace{6mm}
25 ноября 2026 г.
\vspace{17mm}
\end{center}
\noindent{\large\color{darkaccent}\textbf{Итог проверки}}\par
\medskip
Все десять заданий проверены по исходным условиям. Существенных математических ошибок и пропусков не обнаружено. Разбор ошибок и дополнительные упражнения на их закрепление не требуются.
\par\medskip
\noindent\textbf{Рекомендация по оформлению.} При переливании жидкости явно указывайте, что её объём сохраняется. Записывайте единицы измерения величин.
\vfill
\begin{center}
{\small\color{textgray} Лёвин Артём Александрович\\[2mm]
эксклюзивно для Николь Саркисьянц}\par
\vspace{13mm}
\begin{tikzpicture}
\draw[accent,line width=.7pt] (0,0) .. controls (2,0.9) and (4,-0.9) .. (6,0) .. controls (8,0.9) and (10,-0.9) .. (12,0);
\end{tikzpicture}
\end{center}
\end{document}
